\SourceLeaf{399}
\[
\begin{gathered}A^{\mathrm n}=\dfrac{b^{\mathrm n+1}}{b'}\surd{\dfrac{\alpha^{\mathrm n}}{\beta^{\mathrm n}}}\cdot\surd{\alpha\beta},\\[9pt]
B^{\mathrm n}=\dfrac{b^{\mathrm n+1}}{b'}\surd{\dfrac{\beta_{\mathrm n}}{\alpha^{\mathrm n}}}\cdot\surd{\alpha\beta},\\[9pt] H=x\end{gathered}
\]
wird. Die Bestimmung eines reellen $\eta$ ist in \textsc{Gauss} Aufzeichnungen durch eine Lücke unvollendet gelassen,
sie ergibt sich aber, wenn man aus $C=\gamma$, und $D=\delta$ denselben Algorithmus wie eben aus $A$ und $B$ bildet, denn dann sind:
\[
\begin{gathered}
C^{\mathrm n}=\dfrac{b^{\mathrm n+1}}{b'}\surd{\dfrac{\gamma^{\mathrm n}}{\delta^{\mathrm n}}}\surd{\gamma\delta},\\[7pt]
D^{\mathrm n}=\dfrac{b^{\mathrm n+1}}{b'}\surd{\dfrac{\delta^{\mathrm n}}{\gamma^{\mathrm n}}}\surd{\gamma\delta}\\
\begin{aligned}\dfrac{H}{k}y^{\frac12}\eta&=\dfrac{D}{b}\cdot\surd{\dfrac{bC}{aD}}\cdot
{}^{4}\!\surd{\dfrac{b'C'}{a'D'}}\\&\quad\cdot{}^{8}\!\surd{\dfrac{b''C''}{a''D''}}\cdot
{}^{16}\!\surd{\dfrac{b'''C'''}{a'''D'''}}.\;.\;.\end{aligned}\end{gathered}
\]
und diese Grössen $C,\allowbreak D$ nähern sich rasch dem Werthe $\frac{k}{b'}\surd{\gamma\delta}$,
während $\gamma^{\mathrm n}$ und $\delta^{\mathrm n}$ zugleich entweder bedeutend wachsen oder abnehmen,
sobald $\frac{x}{k}y^{\frac12}\eta$ sich von der Einheit unterscheidet.]

\AGMArticle{19}

[Die Beziehungen zwischen den Differentialen der zu untersuchenden Grössen sind besonders einfach,
wenn $a$ und $b$ ungeändert bleiben; es ergibt sich dann unmittelbar aus den Bedingungsgleichungen
zwischen Gliedern mit gleichem Index, dass der Ausdruck
$\frac{1}{2^n}\frac{1}{b^{\mathrm n}}\surd{\frac{\beta^{\mathrm n}\delta_{\mathrm n}}{\alpha^{\mathrm n}\gamma^{\mathrm n}}}\cdot\mathrm{d}\log\frac{\delta^{\mathrm n}}{\beta^{\mathrm n}}$
seinen Werth nicht ändert, wenn darin der Reihe nach $\alpha,\allowbreak \beta,\allowbreak \gamma,\allowbreak \delta$ statt $\beta,\allowbreak \gamma,\allowbreak \delta,\allowbreak \alpha$ gesetzt wird.
Die beiden so erhaltenen Ausdrücke können zu einem dem erstern entsprechenden Ausdruck für den Index $n+1$ vereinigt werden,
der Werth dieses Ausdrucks ist also auch unabhängig vom Index $n$.
Durch Übergang zur Grenze und durch Vergleichung mit den Differentialen der auf verschiedene Weise gebildeten Quotienten ergibt sich
\[
\begin{aligned}\dfrac{1}{k}\mathrm{d}\log\eta
&=\dfrac{1}{a}\surd{\dfrac{\alpha\delta}{\beta\gamma}}\cdot\mathrm{d}\log\dfrac{\delta}{\alpha}
\\[7pt]&=\dfrac{1}{b}\surd{\dfrac{\beta\delta}{\alpha\gamma}}\cdot\mathrm{d}\log\dfrac{\delta}{\beta}
\\[7pt]&=\dfrac{1}{c}\surd{\dfrac{\gamma\delta}{\alpha\beta}}\cdot\mathrm{d}\log\dfrac{\delta}{\gamma}
\\[7pt]&=\dfrac{1}{c}\surd{\dfrac{\delta'}{\beta'}}\cdot\mathrm{d}\log\dfrac{\delta\delta}{\delta'}
\\[7pt]&=\dfrac{1}{a}\surd{\dfrac{\beta\gamma}{\alpha\delta}}\cdot\mathrm{d}\log\dfrac{\gamma}{\beta}
\\[7pt]&=\dfrac{1}{b}\surd{\dfrac{\alpha\gamma}{\beta\delta}}\cdot\mathrm{d}\log\dfrac{\gamma}{\alpha}
\\[7pt]&=\dfrac{1}{c}\surd{\dfrac{\alpha\beta}{\gamma\delta}}\cdot\mathrm{d}\log\dfrac{\alpha}{\beta}
\\[7pt]&=\dfrac{1}{c}\surd{\dfrac{\beta'}{\delta'}}\cdot\mathrm{d}\log\dfrac{\beta'}{\beta\beta}\end{aligned}
\]
eine Relation, welche also immer gilt, wenn die Indices der $a,\allowbreak b,\allowbreak c,\allowbreak \alpha,\allowbreak \beta,\allowbreak \gamma,\allowbreak \delta,\allowbreak \eta$
um gleich viel Einheiten vermehrt werden. Hieraus folgt der von \textsc{Gauss} aufgezeichnete Satz:]
\[
\begin{aligned}\int\dfrac{\mathrm{d}U}{\surd{(aa\sin U^2+bb\cos U^2)}}
&=\int\dfrac{\mathrm{d}V}{\surd{(aa\cos V^2+bb\sin V^2)}}\\[6pt]&=\dfrac{u}{k}\end{aligned}
\]

\SourceLeaf{400}
[Mit Hülfe der Gleichung für $\mathrm{d}\log\frac{\beta'}{\beta\beta}$ lassen sich durch wiederholte Anwendung
derselben unmittelbar die Reihen für die Differentiale der einzelnen Grössen $\alpha,\allowbreak \beta,\allowbreak \gamma,\allowbreak \delta$ aufstellen,
für ein negatives $\delta$ entsteht:
\[
\begin{aligned}\dfrac{k}{2}\dfrac{\mathrm{d}\log\dfrac{\alpha}{x}}{\mathrm{d}u}+2c'\sin 2V'
&=\dfrac{k}{2}\dfrac{\mathrm{d}\log\dfrac{\beta}{x}}{\mathrm{d}u}
\\[7pt]&=\dfrac{k}{2}\dfrac{\mathrm{d}\log\dfrac{\gamma}{x}}{\mathrm{d}u}+a\dfrac{\sin V}{\cos U}
\\[7pt]&=\dfrac{k}{2}\dfrac{\mathrm{d}\log\dfrac{\delta}{x}}{\mathrm{d}u}-b\dfrac{\cos V}{\sin U}
\\[7pt]&=c'\sin 2V'+c''\sin 2^2V''+c'''\sin 2^3V'''+.\;.\;.\text{]}\end{aligned}
\]

\AGMArticle{20}

[Die Differentiale zweiter Ordnung, welche auf die Grösse $\eta$ als einzige unabhängig Veränderliche sich beziehen,
können unmittelbar durch Differentiation der Differentialgleichung erster Ordnung hergeleitet
und die Bestimmung ihrer Werthe in der Weise dargestellt werden, dass die Ausdrücke von der Form
\[
\dfrac{\mathrm{d}}{\mathrm{d}\log\eta}
\left[\dfrac{\mathrm{d}\log\dfrac{\alpha}{\beta}}{\mathrm{d}\log\eta}\right]
+\tfrac12\dfrac{\mathrm{d}\log\dfrac{\alpha\beta}{\gamma\delta}}{\mathrm{d}\log\eta}
\cdot\dfrac{\mathrm{d}\log\dfrac{\alpha}{\beta}}{\mathrm{d}\log\eta}
\]
für die mit irgend welchem gemeinsamen Index behafteten $\alpha,\allowbreak \beta,\allowbreak \gamma,\allowbreak \delta$ verschwinden,
in welcher Reihenfolge die $\alpha,\allowbreak \beta,\allowbreak \gamma,\allowbreak \delta$ auch darin eingesetzt sein mögen.
Multiplicirt man Zähler und Nenner unter dem zweimal zu differentiirenden Logarithmus mit $\beta$,
ersetzt $\alpha\beta$ durch $b'\beta'$ und die Derivirten erster Ordnung durch ihre Werthe
so erhält man für zwei aufeinander folgende Indices:
\[
\dfrac{\mathrm{dd}\log\dfrac{\beta'}{\beta\beta}}{(\mathrm{d}\log\eta)^2}
+2\dfrac{a'c'}{kk}\dfrac{\delta'}{\beta'}
-2\dfrac{ac}{kk}\dfrac{\delta}{\beta}-\tfrac12\dfrac{cc}{kk}=0
\]
und hieraus durch wiederholte Anwendung dieser Gleichung und mit Zuhülfenahme
der in Art. 13 für $\mathrm{d}\log\mathrm{M}(a,\allowbreak b)$ gefundenen Reihe die Gleichung
\[
\begin{aligned}\dfrac{\mathrm{d}\log\dfrac{b}{k}}{\mathrm{d}\log y}
&=\dfrac{\mathrm{dd}\log\dfrac{\alpha}{x}}{(\mathrm{d}\log\eta)^2}-\tfrac12\dfrac{ac}{kk}\dfrac{\gamma}{\alpha}
\\[7pt]&=\dfrac{\mathrm{dd}\log\dfrac{\beta}{x}}{(\mathrm{d}\log\eta)^2}+\tfrac12\dfrac{ac}{kk}\dfrac{\delta}{\beta}
\\[7pt]&=\dfrac{\mathrm{dd}\log\dfrac{\gamma}{x}}{(\mathrm{d}\log\eta)^2}-\tfrac12\dfrac{ac}{kk}\dfrac{\alpha}{\gamma}
\\[7pt]&=\dfrac{\mathrm{dd}\log\dfrac{\delta}{x}}{(\mathrm{d}\log\eta)^2}+\tfrac12\dfrac{ac}{kk}\dfrac{\beta}{\delta}\end{aligned}
\]
welche ihre Gültigkeit behält, wenn $a,\allowbreak b,\allowbreak c,\allowbreak \alpha,\allowbreak \beta,\allowbreak \gamma,\allowbreak \delta,\allowbreak x,\allowbreak y,\allowbreak \eta$
mit irgend einem gemeinsamen von Null verschiedenen Index behaftet werden.

\SourceLeaf{401}
Durch die Vereinigung dieses Resultats mit der im vorigen Artikel gefundenen Entwicklung von
$\mathrm{d}\log\frac{\alpha}{x}$ ergibt sich die von \textsc{Gauss} gefundene Werthausmittelung des Integrals:]
\[
\begin{aligned}
\int\surd{(aa\sin U^2+bb\cos U^2)}\,\mathrm{d}U
\\[7pt]&=\int\dfrac{aabb\,\mathrm{d}V}{(aa\cos V^2+bb\sin V^2)^{\frac32}}\\
&=\dfrac{u}{k}\{a'a'-2c''c''-4c'''c'''-8c''''c''''-.\;.\}\\
&\quad+c'\sin 2V'-c''\sin 4V''\\&\quad-c'''\sin 8V'''-c''''\sin 16V''''-.\;.\;.\end{aligned}
\]

\AGMArticle{21}

[Bildet man die Gleichungen zwischen den Derivirten der $\alpha,\allowbreak \beta,\allowbreak \gamma,\allowbreak \delta$
nach der Grösse $y$ als unabhängig Veränderliche, so müssen dieselben Coëfficienten dieser Derivirten entstehen,
wie in den Gleichungen für die Differentiale nach $\eta$ und da die letztern Gleichungen in solche Form gebracht werden können,
dass die Verhältnisse zwischen den Derivirten von Quotienten bestimmt werden,
so müssen die erwähnten Coëfficienten auch diesen Derivirten umgekehrt proportional sein.
Ersetzt man sie durch deren reciproken Werthe, so ersieht man, dass auch die übrigen Glieder jener Gleichungen
durch solche Derivirten dargestellt werden können und zwar so, dass der Ausdruck von der Form
\[
\dfrac{\left[\dfrac{\mathrm{d}\log\dfrac{\delta}{\beta}}{\mathrm{d}\log y}\right]}
{\left[\dfrac{\mathrm{d}\log\dfrac{\delta}{\beta}}{\mathrm{d}\log\eta}\right]}
-\tfrac12\dfrac{\mathrm{d}\log\dfrac{\alpha\beta}{xx}}{\mathrm{d}\log\eta}
\]
für einen beliebigen gemeinsamen Index der Zeichen $\alpha,\allowbreak \beta,\allowbreak \gamma,\allowbreak \delta,\allowbreak x,\allowbreak y,\allowbreak \eta$
seinen Werth nicht ändert, in welcher Reihenfolge man auch die Grössen $\alpha,\allowbreak \beta,\allowbreak \gamma,\allowbreak \delta$ mit einander vertauscht.
Vergleicht man diesen Ausdruck, welcher sich auf $\frac{\delta}{\beta},\allowbreak \frac{\alpha\gamma}{xx},\allowbreak \eta,\allowbreak y$ bezieht,
mit demjenigen für $\frac{\delta'}{\beta'}=\frac{\gamma\delta}{\alpha\beta},\allowbreak \frac{\alpha'\gamma'}{x'x'},\allowbreak \eta'=\eta\eta,\allowbreak y'=yy$
und beachtet, dass nach Art. 19
\[
\begin{gathered}\left[\dfrac{\mathrm{d}\log\dfrac{\alpha}{\beta}}{\mathrm{d}\log\eta}\right]^2
=\dfrac{\mathrm{d}\log\dfrac{\gamma'}{\beta'}}{\mathrm{d}\log\eta}\cdot
\dfrac{\mathrm{d}\log\dfrac{\alpha'}{\beta'}}{\mathrm{d}\log\eta},\\[9pt]
\left[\dfrac{\mathrm{d}\log\dfrac{\delta}{\gamma}}{\mathrm{d}\log\eta}\right]^2
=\dfrac{\mathrm{d}\log\dfrac{\delta'}{\alpha'}}{\mathrm{d}\log\eta}\cdot
\dfrac{\mathrm{d}\log\dfrac{\delta'}{\gamma'}}{\mathrm{d}\log\eta}\end{gathered}
\]
ist, so ergibt sich, dass der Werth jenes Ausdrucks auch unabhängig von dem Index der Zeichen
$\alpha,\allowbreak \beta,\allowbreak \gamma,\allowbreak \delta,\allowbreak x,\allowbreak y,\allowbreak \eta$ und, wie aus den Grenzwerthen folgt, gleich Null sein muss.
